\documentclass[withoutpreface,bwprint]{cumcmthesis}

\usepackage{amsmath,amssymb}
\usepackage{graphicx}
\graphicspath{{figures/}{./}}
\usepackage{booktabs}
\usepackage{multirow}
\usepackage{longtable}
\usepackage{tabularx}
\usepackage{float}
\usepackage{subcaption}
\usepackage{url}
\usepackage{listings}
\hypersetup{hidelinks}

\title{面向算电协同的多目标调度优化研究}
\author{参赛队号：\quad 题号：C}
\date{}

\begin{document}

\maketitle

\begin{abstract}
针对六区域异构 AI 任务与新能源、储能及电网协同调度问题，本文建立了以任务执行区域、开工时刻和储能动作为决策变量，以运行成本、碳排放、网络时延、新能源利用率、服务质量及峰值净购电为评价指标的多目标模型，并纳入算力、功率、时延、期限、能量平衡和 SOC 边界约束。求解阶段采用固定 Huber 滞后回归、确定性贪心任务调度、规则型 SOC 前向策略和固定种子情景复算，所得结果仅为通过硬约束审计的启发式可行方案。结果表明，$50000$ 条任务记录均通过输入校验，缺失单元格数量为 $0$，预测总体 MAE 和 RMSE 分别为 $24.9733$~GPU 和 $35.6945$~GPU。问题二所得运行成本为 $-419658147.7076$~元，碳排放为 $14.3916$~tCO$_2$，网络时延为 $5.1594$~ms，新能源利用率为 $0.6791$，迁移任务数为 $209$ 个。问题三累计充电 $2076.3155$~MWh、放电 $0$~MWh，含储能方案的运行成本增加 $304223.9291$~元，碳排放与峰值净购电变化均为 $0$。问题四联合基准的运行成本、碳排放、新能源利用率和峰值净购电分别为 $-419353923.7785$~元、$14.3916$~tCO$_2$、$0.6791$ 和 $40.4664$~MW，并实际检查了 $15$ 个情景。研究结果说明，该方法可生成满足已审计硬约束的联合调度方案，但其属于启发式复算结果，不构成 Pareto 解，且不具有最优性证书。
\end{abstract}

\noindent\textbf{关键词：}异构 AI 任务；算力调度；Huber 回归；储能协同；碳感知调度；情景分析

\clearpage

\section{问题重述}

随着人工智能推理与训练任务持续增长，数据中心的算力分配已经与能源供给、碳排放和通信服务质量形成紧密耦合。题目给出了六个区域的 GPU 容量、IT 与设施功率边界、逐时非 AI 负荷、PUE、电价、碳强度、新能源出力、购售电限制、储能参数以及区域间单向网络时延，同时提供了三类异构 AI 任务的到达时刻、GPU 需求、预计持续时间、来源区域、最大允许时延和完成期限。需要在满足物理边界与服务约束的条件下，合理确定任务的执行区域和开工时刻，并协调新能源、储能及电网购售电行为。

整个运行过程具有明确的分层时间边界。第 $0$ 至第 $2399$ 小时构成任务到达主时域，第 $2400$ 至第 $2405$ 小时不再接收新任务，仅允许此前到达的弹性任务继续执行并完成结清；第 $2406$ 小时不得安排任何任务，只用于电力运行和储能终端状态结算。所有任务均不可拆分、不可抢占，也不得在执行过程中迁移。实时推理任务必须在到达时刻立即开工，批量推理与 AI 训练任务可以在允许范围内延迟，但其完成时点不得超过任务截止时间和系统统一终点。

对于问题一，需要按区域和任务类型统计 GPU 需求的规模、分布及时间规律，建立短期需求预测模型，并按照时间顺序完成模型训练、验证与最终测试。预测结果仅用于模型精度评价，最后 $24$ 小时的基础调度必须使用第 $2376$ 至第 $2399$ 小时的实际到达任务。调度时还需根据分钟级持续时间计算任务与各小时的真实重叠，据此分别审计瞬时 GPU 占用和小时电量，并展示任务实际起止时刻、执行区域及区域 GPU 利用情况。

对于问题二，需要使用全体实际任务和逐时电力参数重新构造碳感知调度方案。在任务唯一执行、实时任务即时开工、网络时延、GPU 容量、IT 功率、设施功率及完成时限均得到满足的基础上，系统运行成本、购电碳排放、有效网络时延和新能源利用率应被统一复算。问题二不配置储能，新能源被用于直接供能、区域外送或弃置，不足负荷由电网购电补足，且购电与售电行为必须服从区域边界及互斥要求。

对于问题三，任务调度结果不再作为决策变量，区域 AI IT 负荷被固定为附件中的基准负荷，区域设施负荷由基准 AI IT 负荷与非 AI IT 负荷之和乘以 PUE 得到。需要在储能初始状态、荷电状态递推、容量和功率边界、充放电互斥、购售电互斥以及终端能量不透支等条件下，确定逐区域逐时的充电、放电、购电、售电、新能源分配和弃电策略，并将含储能方案与无储能比较状态在成本、碳排放、峰值净购电功率和负荷爬坡方面进行比较。

对于问题四，需要把任务迁移与开工、设施负荷、新能源分配、储能运行以及购售电决策纳入同一框架，在运行成本、碳排放、网络时延、服务质量、新能源利用率和区域峰值净购电之间进行权衡。碳约束、电价机制与新能源波动等情景应从同一联合基准独立重构，并在相同任务集、容量边界、功率映射和评价口径下接受完整审计，从而辨析不同外部条件对任务迁移方向、执行窗口与能源行为的影响。

\section{问题分析}

本题的本质是一个跨时间、跨区域的算力与能源耦合调度问题。任务侧包含执行区域和开工时刻等离散选择，能源侧包含新能源分配、购售电和储能功率等连续决策，而任务不可抢占、瞬时资源容量、分钟级执行时长与小时级能源结算又使两侧具有不同的计量口径。因此，建模时不能直接以逐小时汇总 GPU 需求替代任务，也不能把所有触及某一小时的任务均视为完整运行一小时。合理的处理方式是先构造每个任务的合法区域--开工候选域，再以任务是否与某小时相交的指示量审计瞬时容量，以实际重叠时长计算 GPU-hour、AI 电量及平均设施负荷，最后将该负荷传递至新能源、储能与电网模型。

\subsection{问题一分析}

问题一同时包含描述性统计、时间序列预测和末端基础调度，但三项工作使用的数据口径并不相同。统计阶段需要把任务到达记录转换为区域--任务类型--小时序列，以任务数量、GPU 需求、GPU-hour、执行时长、到达强度、峰值并发及周期相关性刻画异构需求。预测阶段应严格采用顺序划分，先利用第 $0$ 至第 $2351$ 小时拟合候选模型，再以第 $2352$ 至第 $2375$ 小时完成模型比较；模型结构确定后，应在规定的历史区间上重新拟合，并只将第 $2376$ 至第 $2399$ 小时作为最终测试区间。由于部分区域和任务类型可能出现零需求，WAPE 只有在真实需求绝对值之和为正时才具有意义，故 MAE、RMSE 与带分母门禁的 WAPE 应共同用于误差评价。

需求序列可能同时呈现日周期、周周期、短期惯性和突发扰动。为此，预测量被表示为区域与任务类型的逐小时到达 GPU 需求，模型中应包含短期滞后、日内滞后、周内滞后、正余弦周期项以及区域和任务类型效应。考虑到尖峰任务会显著影响平方损失，固定 Huber 回归可被用于降低异常峰值对参数估计的支配作用。预测输出在模型拟合后才执行非负截断，以避免截断操作改变训练目标。实际实现采用固定 Huber 滞后回归，未执行完整候选网格选模，因此最终测试误差只能说明该已实现模型在保留测试时段上的表现。

最后 $24$ 小时的调度不能由预测需求替代，而应直接读取同期实际到达任务。对于任务 $i$，其持续小时数由分钟时长换算得到，有效截止时点取任务有效截止与系统终点中的较早者；随后根据到达时刻、任务类型、单向时延和完成边界生成合法执行区域与整数开工时刻。实时推理任务的开工时刻被固定为到达时刻，弹性任务则可在候选域内延迟执行。实际算法采用确定性贪心启发式，根据等待时间、网络时延、峰值 GPU 利用率及固定词典序选择候选，并通过完整资源审计构造可行方案；契约中规划的有限冲突回溯并未实现，故若贪心构造失败，只能报告在既定方法内未能确认可行性。

资源约束需要区分瞬时口径和能量口径。只要任务在小时 $t$ 内存在正重叠，其持续占用的 GPU 数及相应瞬时 IT 功率就应被计入容量审计；任务在该小时的实际重叠时长则用于形成 GPU-hour 和 AI 电量。区域总 IT 功率由 AI 瞬时功率与不可迁移的非 AI IT 负荷相加，设施功率由同一总 IT 功率乘以区域 PUE 得到。甘特图应以任务实际开工和完成时点绘制，并同时给出峰值 GPU 利用率与时间加权平均利用率，从而避免仅以小时平均值掩盖短时容量冲突。

\subsection{问题二分析}

问题二需要以第 $0$ 至第 $2399$ 小时的实际任务重新进行全时域调度，并将第 $2400$ 至第 $2405$ 小时发生的末端任务能耗纳入成本和碳排放核算。任务侧沿用问题一的合法候选域、唯一执行、实时即时开工、截止时间、单向时延及瞬时容量约束，但不能直接沿用预测结果。AI 任务的小时电量由 GPU 需求、任务类型功率映射与真实重叠时长共同计算，区域设施平均负荷则由 AI 平均 IT 负荷和固定非 AI IT 负荷之和乘以 PUE 获得。

由于问题二不配置储能，逐时能源分配可在给定设施负荷后被确定性复算。新能源先被用于满足本地负荷；当新能源不足时，缺口由电网购电补足；当新能源富余时，可在区域外送上限内售出，其余部分被记为弃电。购电与售电不能在同一时段同时发生，超过进口上限的任务候选应被判定为不可行。系统成本按照购电支出扣除售电收入计算，碳排放只由电网购电量与对应碳强度相乘累积，新能源利用率的分子包含直接使用、允许计入的新能源充电和新能源外送，而分母为可用新能源总量。

成本、碳排放、网络时延、新能源未利用程度与等待量具有不同量纲，因而不能直接相加。方法契约给出的处理是以相对基准真实差值的有限正中位数作为尺度，对有效维度进行标准化后加权评价；若某一指标不存在正尺度或其定义分母为零，该维度应退出评分。网络时延采用 GPU 需求与执行时长加权的单向时延，只有权重和为正时才计算。实际实现并未执行规定数量的单任务迁移或开工邻域改进，而是对确定性任务方案进行完整能源复算，因此所得方案属于启发式可行方案，不能被解释为连续参数或离散候选空间中的最优解。

\subsection{问题三分析}

问题三与前两问的关键区别在于任务负荷已经固定。区域 IT 负荷应直接取基准 AI IT 负荷与非 AI IT 负荷之和，设施负荷再乘以区域 PUE，任务迁移和开工时刻不得被重新优化。这样，问题三可被转化为逐区域、逐小时的能源与储能协同问题。无储能状态只用于建立比较基准，即使该状态违反购电边界，也不能据此截断含储能方案的搜索；含储能方案必须从附件给定的绝对初始 SOC 独立向前构造。

能源侧需要同时满足新能源守恒、负荷平衡和储能动态。可用新能源被拆分为直接使用、储能充电、外送和弃电，电网购电则被拆分为直接供负荷和储能充电。储能总充电功率等于新能源充电与电网充电之和，时段末 SOC 由前一状态、充电效率和放电效率递推。每次动作后还需检查未来剩余最大充电能力能否保证终端 SOC 不低于初始 SOC，以避免前期过度放电造成末端不可恢复。充电、放电和空闲三种状态相互排斥，购电与售电也应通过状态变量或等价逻辑保持互斥。

含储能与无储能状态的比较需要统一指标口径。净购电被定义为购电减售电，区域峰值指标取整个时域净购电的非负最大值，负荷波动采用相邻小时净购电绝对变化量之和，等效循环次数仅在储能容量为正时计算。运行成本、碳排放和新能源利用率均应由各自方案的能源流重新计算。实际实现采用规则型 SOC 前向策略，并未执行完整动作邻域搜索，因而储能策略只能被表述为满足已审计边界的规则型候选；若含储能成本高于无储能状态，也不能将其包装为经济性改进。

\subsection{问题四分析}

问题四把前三问中的任务调度、能源平衡与储能动态统一起来。任务区域和开工时刻决定 AI 小时电量及设施负荷，设施负荷进一步决定新能源消纳、电网购电与储能动作，能源价格和碳强度又会反向影响弹性任务的迁移和延迟选择。由于离散任务决策与连续能源决策跨越完整时域，直接穷举全部组合并不可行，因此应先删除违反时延、到达、截止和资源下界的候选，再独立构造全时域联合可行初始方案，并通过有界交替邻域对任务与储能动作进行改进。

联合评价包含运行成本、购电碳排放、有效网络时延、有效等待损失、新能源未利用率及六区域非负峰值净购电功率。服务质量中的实时即时开工率、SLA 满足率和按期完成率本质上属于硬约束审计，只有分母为正且候选之间存在差异的等待分量才适合作为评分项。实际实现以平均等待构造有限代理，并未实现完整等待质量公式。各指标应采用有限正差值尺度进行无量纲化，分母无效或候选间无变化的指标不得被强行纳入综合评分。

情景比较需要保持单因素可比。不同碳约束、电价机制和新能源水平或波动情景均应从同一联合基准独立重建，在相同任务集、容量、时延、PUE 和功率映射下复算全部指标。随机新能源波动使用固定种子，以保证重复运行的一致性。实际方法未执行约定的共享交替邻域搜索，每个情景只重建并审计一个规则型候选，且问题四的任务迁移方案保持联合基准不变。因此，未达到某一碳目标只能说明声明的启发式候选未满足该目标，不能据此推断该情景在数学上不可行；情景结果也只构成规则型策略比较，不构成 Pareto 前沿。

\subsection{子问题间的衔接关系}

四个子问题共享任务重叠计算、功率换算与约束审计口径，但其决策变量和输入负荷不能相互混用。问题一的预测模型可独立训练，预测输出只服务于精度评价，最后 $24$ 小时基础调度仍由实际任务生成。问题二可以复用问题一的候选域与容量审计方法，但必须以全时域实际任务和逐时能源数据重新计算。问题三固定采用附件中的基准 AI IT 负荷，因而可与任务预测和迁移调度相对独立地求解。问题四则继承任务约束、能源守恒、储能状态和统一指标定义，但联合方案必须独立构造，不能简单拼接问题二的任务方案与问题三的储能方案。

\section{模型假设}

为统一小时级调度与分钟级资源计量，假设任务开工时刻按照整数小时选取，而任务在各小时内的 GPU-hour 和能耗根据实际起止时刻与该小时的重叠时长计算。附件中的 GPU 需求、预计持续时间、到达时刻、来源区域、最大允许时延和最晚完成时刻均被视为准确且确定的输入参数，不另行建立执行时间和资源需求的随机模型。所有任务仅能选择一个执行区域和一个开工时刻，自开工至完成保持连续运行，不允许拆分、抢占、暂停或中途迁移。

假设实时推理任务必须在到达时刻立即开工，其执行区域只能从满足单向网络时延和资源容量的区域中选择。批量推理与 AI 训练任务可以延迟开工，但不得早于到达时刻，且必须在任务有效截止时点与系统时点 $2406$ 中较早者之前完成。第 $2400$ 至第 $2405$ 小时只用于此前到达弹性任务的结清，第 $2406$ 小时不产生任务资源占用，只承担电力与储能终端结算。

假设同一任务类型的单位等效 GPU 平均 IT 功率在研究期内保持为功率映射附件给定值，不随区域和时间改变。非 AI IT 负荷在各区域和各时段固定且不可迁移，只有 AI 任务负荷能够随执行区域及开工时刻改变。区域设施功率由 AI 与非 AI 合计 IT 功率乘以对应 PUE 得到，PUE 在研究期内按附件取值，不建立其随环境或负荷率变化的动态模型。

假设区域间通信只体现为附件给出的来源区域至执行区域单向网络时延，时延在研究期内确定且不受网络拥塞影响。来源区域与执行区域相同时仍采用时延矩阵中的对应记录，不自行设定为零。由于题面没有给出网络带宽、迁移数据量、传输能耗和传输费用等信息，上述因素不被纳入模型。

假设购电价格、售电价格、电网碳强度、可用新能源和固定负荷在单个小时内保持恒定。各区域逐小时独立满足能源平衡，不建立区域间线路潮流模型，区域售电被视为向系统边界外送。可用新能源可以被分配至直接供给负荷、储能充电、外送或弃电，其优先关系由实际策略与能源平衡共同确定，而不预设全部新能源均被本地直接消纳。运行碳排放仅按电网购电量与相应碳强度累计，不计新能源、储能设备和通信设施的生命周期排放。

假设储能在第 $0$ 小时运行前的状态唯一取自附件中的初始 SOC，后续状态均在该绝对值基础上递推，不将初始状态归零或作为待估参数。区域逐时基准 SOC 只用于方案核验与比较，不作为优化轨迹的固定约束。充电效率作用于进入储能的能量，放电效率作用于储能侧能量消耗，未给出的自放电与储能退化成本不予考虑。每个区域在第 $2406$ 小时运行后的 SOC 不低于其初始 SOC，从而避免通过透支终端能量获得虚假的当期收益。

假设同一储能在同一小时只能处于充电、放电或空闲状态，不能同时充放电；区域购电和售电也不能在同一小时同时发生。两类行为均受附件所给功率上限约束，若价格和效率条件不能自然排除无效循环，则通过互斥状态约束加以限制。问题三中的 IT 负荷固定为基准 AI IT 负荷与非 AI IT 负荷之和，不重新调整任务的执行区域和开工时刻。

假设预测模型与正式调度严格分离，预测区间的估计值只用于测试精度评价，最后 $24$ 小时基础调度以及问题二至问题四的正式计算均使用实际任务和实际逐时参数。服务质量仅由题内可观测的即时开工、单向时延 SLA、按期完成情况及弹性任务等待时间表示，不引入用户满意度或可靠性等缺乏数据支持的变量。进行碳约束、电价机制与新能源波动情景比较时，除指定变化因素外，其余任务集、设备容量、网络时延、功率映射及评价公式均保持一致。

假设普通空白值在未经核验前均属于不可裁决缺失，不自动替换为零。只有在互斥进流与出流字段中，空白所在行为孤立事件、对应另一列明确非零，且备注、相邻记录和全表口径共同支持互斥解释时，才将其认定为结构零。经过处理的数据还应接受任务唯一执行、实时即时开工、时延与截止、GPU 和功率容量、能源守恒、SOC 递推、终端状态及购售电边界的统一审计。

\section{符号说明}

本文使用的主要符号如表~\ref{tab:symbols} 所示。对于同时具有输入参数和派生量的符号，其具体属性在表中一并注明；功率变量均以区域、小时为索引，能量累计时采用时段长度 $\Delta t$ 进行换算。

\small
\begin{longtable}{lll}
\caption{主要符号及其含义}
\label{tab:symbols}\\
\toprule
符号 & 含义 & 类型或单位 \\
\midrule
\endfirsthead
\toprule
符号 & 含义 & 类型或单位 \\
\midrule
\endhead
\bottomrule
\endfoot
$i,r,t,s,k$ & 任务、区域、小时、开工时刻和任务类型索引 & 索引 \\
$\mathcal R$ & 六个计算区域构成的集合 & 集合 \\
$\mathcal K$ & 三类异构 AI 任务构成的集合 & 集合 \\
$\mathcal T_A$ & 任务到达主时域 & $\{0,\ldots,2399\}$ \\
$\mathcal T_C$ & 末端任务结清时域 & $\{2400,\ldots,2405\}$ \\
$\mathcal T_E$ & 电力与储能结算时域 & $\{0,\ldots,2406\}$ \\
$\mathcal S$ & 合法整数开工时刻集合 & $\{0,\ldots,2405\}$ \\
$\Delta t$ & 单个电力结算时段长度 & h \\
$a_i$ & 任务 $i$ 的到达小时 & h \\
$d_i,h_i$ & 任务持续分钟数及换算后的持续小时数 & min，h \\
$g_i$ & 任务 $i$ 持续占用的等效 GPU 数 & GPU \\
$k_i,o_i$ & 任务 $i$ 的类型和来源区域 & 类别，区域 \\
$F_i$ & 任务 $i$ 的有效完成截止时点 & h \\
$W_i^{\max}$ & 任务 $i$ 的最大允许等待时间 & h \\
$\ell_i^{\max}$ & 任务 $i$ 允许的最大网络时延 & ms \\
$\ell_{or}$ & 来源区域 $o$ 至执行区域 $r$ 的单向时延 & ms \\
$\omega_{it}(s)$ & 任务 $i$ 在开工时刻 $s$ 下与小时 $t$ 的重叠时长 & h \\
$\chi_{it}(s)$ & 任务 $i$ 在小时 $t$ 是否存在瞬时占用 & 二元参数 \\
$\Omega_i$ & 任务 $i$ 的合法区域--开工候选域 & 有限集合 \\
$x_{irs}$ & 任务 $i$ 是否在区域 $r$ 于时刻 $s$ 开工 & 二元变量 \\
$s_i,c_i,r_i$ & 任务实际开工时刻、完成时刻和执行区域 & h，h，区域 \\
$D_{rkt}$ & 区域 $r$ 类型 $k$ 在小时 $t$ 的实际到达 GPU 需求 & GPU \\
$D'_{rkt}$ & 用于滞后回归的历史需求输入 & GPU \\
$\widehat D_{rkt}$ & 区域--任务类型逐小时 GPU 需求预测值 & GPU \\
$u_{rkt}$ & Huber 回归截断前的线性预测量 & GPU \\
$z_t$ & 小时 $t$ 的日周期与周周期特征向量 & 无量纲 \\
$\beta,\gamma,\alpha_r,\delta_k$ & 回归系数、周期系数、区域效应和类型效应 & 参数 \\
$\lambda$ & Huber 回归的正则化系数 & 无量纲 \\
$\sigma_{rk}$ & 区域--任务类型序列的残差尺度 & GPU \\
$\rho_{\delta}(\cdot)$ & 阈值为 $\delta$ 的 Huber 损失函数 & 无量纲 \\
$G_r^{\max}$ & 区域 $r$ 的可调度 GPU 容量 & GPU \\
$P_r^{IT,\max}$ & 区域 $r$ 的最大 IT 功率 & MW \\
$P_r^{F,\max}$ & 区域 $r$ 的最大设施功率 & MW \\
$p_k^{GPU}$ & 类型 $k$ 的单位等效 GPU 平均 IT 功率 & MW/GPU \\
$\pi_r$ & 区域 $r$ 的能源使用效率 PUE & 无量纲 \\
$L_{rt}^{N}$ & 区域 $r$ 的不可迁移非 AI IT 负荷 & MW \\
$L_{rt}^{BAI}$ & 问题三固定的基准 AI IT 负荷 & MW \\
$L_{rt}^{AI,cap}$ & AI 任务形成的瞬时 IT 功率 & MW \\
$L_{rt}^{AI,avg}$ & AI 任务形成的小时平均 IT 负荷 & MW \\
$E_{rt}^{AI}$ & AI 任务在小时 $t$ 消耗的电量 & MWh \\
$L_{rt}^{F,avg}$ & 区域设施侧小时平均负荷 & MW \\
$U_{rt}^{peak}$ & 区域小时峰值 GPU 利用率 & 无量纲 \\
$U_{rt}^{avg}$ & 区域小时 GPU-hour 平均利用率 & 无量纲 \\
$R_{rt}^{av}$ & 区域可用新能源功率 & MW \\
$R_{rt}^{use}$ & 直接满足设施负荷的新能源功率 & MW \\
$R_{rt}^{ch}$ & 用于储能充电的新能源功率 & MW \\
$R_{rt}^{sell}$ & 新能源外送功率 & MW \\
$R_{rt}^{curt}$ & 弃置的新能源功率 & MW \\
$G_{rt}^{load}$ & 电网直接供给设施负荷的功率 & MW \\
$G_{rt}^{ch}$ & 电网供给储能充电的功率 & MW \\
$G_{rt}^{buy}$ & 区域从电网购入的总功率 & MW \\
$G_{rt}^{sell}$ & 区域向系统边界外送的功率 & MW \\
$I_r^{\max}$ & 区域最大购电功率 & MW \\
$X_r^{grid}$ & 区域电网侧最大外送功率 & MW \\
$X_r^{storage}$ & 储能信息限定的售电功率上限 & MW \\
$y_{rt}^{grid}$ & 区域购电与售电互斥状态 & 二元变量 \\
$C_{rt}$ & 储能总充电功率 & MW \\
$D_{rt}$ & 储能放电功率 & MW \\
$S_{rt}$ & 区域储能在小时 $t$ 运行后的 SOC & MWh \\
$B_r$ & 区域储能额定容量 & MWh \\
$S_r^0$ & 第 $0$ 小时运行前的初始 SOC & MWh \\
$S_r^{min},S_r^{max}$ & 区域储能 SOC 下限与上限 & MWh \\
$C_r^{\max},D_r^{\max}$ & 区域最大充电与放电功率 & MW \\
$\eta_r^c,\eta_r^d$ & 区域储能充电效率与放电效率 & 无量纲 \\
$N_{rt}$ & 区域净购电功率 $G_{rt}^{buy}-G_{rt}^{sell}$ & MW \\
$P_r^{peak}$ & 区域时域内非负峰值净购电功率 & MW \\
$V_r$ & 区域净购电绝对爬坡量 & MW \\
$C_{\mathrm{op}}$ & 购电支出扣除售电收入后的运行成本 & 元 \\
$E_{\mathrm{CO2}}$ & 电网购电产生的累计碳排放 & tCO$_2$ \\
$L_{\mathrm{net}}$ & 按 GPU 需求与执行时长加权的网络时延 & ms \\
$W_{\mathrm{net}}$ & 网络时延指标的有效权重和 & GPU$\cdot$h \\
$U_{\mathrm{RE}}$ & 新能源利用率 & 无量纲 \\
$W_{\mathrm{RE}}$ & 新能源利用率的可用新能源总量 & MWh \\
$Q_{\mathrm{RT}}$ & 实时任务即时开工满足率 & 无量纲 \\
$Q_{\mathrm{SLA}}$ & 网络时延 SLA 满足率 & 无量纲 \\
$Q_{\mathrm{deadline}}$ & 任务按期完成率 & 无量纲 \\
$Q_{\mathrm{wait}}$ & 弹性任务等待质量指标 & 无量纲 \\
$Q_{\mathrm{service}}$ & 综合服务质量指标 & 无量纲 \\
$\Psi_m$ & 候选方案 $m$ 的标准化综合评分 & 无量纲 \\
$w_j$ & 第 $j$ 个有效评价维度的权重 & 无量纲 \\
$d_j$ & 第 $j$ 个评价维度的有限正尺度 & 对应指标单位 \\
$a_{jm}$ & 候选方案 $m$ 在指标 $j$ 上的取值 & 对应指标单位 \\
$m_0$ & 标准化评分使用的基准方案索引 & 索引 \\
$N_{\mathrm{back}}$ & 冲突回溯节点预算 & 节点 \\
$N_{\mathrm{move}}^{q2}$ & 问题二单任务候选评价预算 & 次 \\
$N_{\mathrm{move}}^{q3}$ & 问题三储能动作候选评价预算 & 次 \\
$N_{\mathrm{move,total}}^{q4}$ & 问题四基准与全部情景共享的评价预算 & 次 \\
$D_{\mathrm{search}}$ & 启发式搜索的内部截止时刻 & 时间点 \\
$\rho_C$ & 情景中的碳约束水平参数 & 无量纲 \\
$\delta_P$ & 峰谷电价差调整系数 & 无量纲 \\
$\delta_S$ & 售电价格调整系数 & 无量纲 \\
$\delta_R$ & 新能源出力水平调整系数 & 无量纲 \\
$\xi$ & 新能源波动情景中的固定种子随机扰动 & 无量纲 \\
$\kappa_{rt}$ & 区域逐时电网购电碳强度 & tCO$_2$/MWh \\
$c_{rt}^{buy},c_{rt}^{sell}$ & 区域逐时购电价格和售电价格 & 元/MWh \\
$\mathrm{EFC}_r$ & 区域储能的等效完整循环次数 & 次 \\
\end{longtable}
\normalsize

表~\ref{tab:symbols} 同时区分了瞬时容量变量与小时能量变量，其中 $\chi_{it}(s)$ 被用于 GPU、IT 功率和设施功率的峰值审计，$\omega_{it}(s)$ 则被用于 GPU-hour 与 AI 电量计算。该区分能够保留非整小时任务的真实资源占用，避免小时平均功率掩盖瞬时超限。

能源变量按来源和去向被进一步拆分，使新能源守恒、电网购电成本和储能充电来源能够分别核算。由此，新能源利用率不会把弃电计入有效利用量，碳排放也不会遗漏电网为储能充电所提供的电量。

评分符号仅作用于分母有效且在候选之间存在真实差异的指标。若 $W_{\mathrm{net}}$、$W_{\mathrm{RE}}$ 或服务质量分量的相应分母不为正，则该指标被标记为不可用并退出综合评分，从而避免无依据的除零处理。

% MMW-METHOD-CONTRACT: ee95080ed721168d72ddba0989ae7188b6e8baaceecbefd9c5b20fd61ea1aadf
\section{模型的建立与求解}

\subsection{数据预处理与统一计量口径}

工作负载、GPU 中心、网络时延、功率映射、区域逐时数据与储能信息首先被按字段说明进行一致性核验。工作负载共包含 $50000$ 条输入任务记录，均通过数据完整性校验，缺失单元格数量为 $0$。区域、任务类型和小时索引被统一编码，任务类型与单位等效 GPU 功率按照功率映射表连接，区域间通信采用有向单向时延矩阵。逐时电力数据覆盖任务主时域、末端结清时域和终端结算时点，其中第 $2406$ 小时不得安排任务，仅用于能源与储能状态结算。

任务 $i$ 的持续小时数及有效截止时点被定义为
\begin{equation}
h_i=\frac{d_i}{60},\qquad F_i=\min\{f_i^{\mathrm{valid}},2406\}.
\label{eq:duration_deadline}
\end{equation}
若任务在整数时刻 $s$ 开工，则其与小时区间 $[t,t+1)$ 的实际重叠时长和瞬时占用指示量分别为
\begin{align}
\omega_{it}(s)&=\max\{0,\min(t+1,s+h_i)-\max(t,s)\},\\
\chi_{it}(s)&=\mathbf 1\{\omega_{it}(s)>0\}.
\end{align}
其中，$\omega_{it}(s)$ 被用于 GPU-hour 和电量计算，$\chi_{it}(s)$ 被用于 GPU、IT 功率与设施功率的瞬时容量审计。

系统成本、购电碳排放、有效网络时延和新能源利用率统一按
\begin{align}
C_{\mathrm{op}}
&=\sum_{r,t}\left(c_{rt}^{buy}G_{rt}^{buy}
-c_{rt}^{sell}G_{rt}^{sell}\right)\Delta t,\\
E_{\mathrm{CO2}}
&=\sum_{r,t}\kappa_{rt}G_{rt}^{buy}\Delta t,\\
L_{\mathrm{net}}
&=\frac{\sum_i g_ih_i\ell_{o_ir_i}}{\sum_i g_ih_i},\\
U_{\mathrm{RE}}
&=\frac{\sum_{r,t}(R_{rt}^{use}+R_{rt}^{ch}+R_{rt}^{sell})\Delta t}
{\sum_{r,t}R_{rt}^{av}\Delta t}
\end{align}
复算。网络时延和新能源利用率仅在分母为正时计算，否则相应指标被标记为不可用并退出评分。多目标候选只有通过全部硬约束后才参与比较，其指标差异采用有限正中位绝对尺度标准化\cite{huber1964robust,miettinen1999nonlinear}。

% MMW-ALGORITHM: 固定Huber滞后回归；50000任务NumPy区间切片确定性贪心调度；按真实调度重建AI小时电量和设施负荷；问题二从真实附件中定位区域级MaxGridImport_MW和MaxGridExport_MW并执行购售电边界及互斥审计；问题三使用Baseline AI负荷执行无储能比较与规则型SOC前向策略；问题四使用调度派生FacilityLoad执行联合储能前推并复算情景；MMW_PILOT只读取真实输入并执行有限样本门禁

实际程序采用 Python 与 scikit-learn 实现。需求预测采用固定 Huber 滞后回归，任务调度采用 NumPy 区间切片支持的确定性贪心构造，能源流根据真实调度形成的小时电量重新计算。问题三采用规则型 SOC 前向策略，问题四依据调度派生设施负荷复算各情景；除新能源波动情景采用固定随机种子 $2026$ 外，其余过程均为确定性构造。

\subsection{问题一：需求预测与基础算力调度}

\subsubsection{GPU 需求统计与预测模型}

% MMW-ID: OBJ-Q1
区域 $r$、任务类型 $k$ 在小时 $t$ 的到达 GPU 需求被聚合为
\begin{equation}
D_{rkt}=\sum_{i:o_i=r,\;k_i=k,\;a_i=t}g_i.
\end{equation}
六区域与三类任务的统计结果如表~\ref{tab:q1_group_statistics} 所示。

\begin{table}[htbp]
\centering
\caption{六区域三类任务的需求统计}
\label{tab:q1_group_statistics}
\begin{tabular}{llrrr}
\toprule
来源区域 & 任务类型 & 任务数 & GPU 需求 & GPU-hour\\
\midrule
RegionA & AITraining & 879 & 63070 & 214319.2833\\
RegionA & BatchInference & 3331 & 45066 & 154341.4500\\
RegionA & RealTimeInference & 5852 & 23585 & 79791.7167\\
RegionB & AITraining & 871 & 62551 & 212294.6333\\
RegionB & BatchInference & 3312 & 44274 & 151260.4333\\
RegionB & RealTimeInference & 5017 & 19989 & 67203.0833\\
RegionC & AITraining & 815 & 58899 & 198331.9667\\
RegionC & BatchInference & 2523 & 34710 & 117897.3667\\
RegionC & RealTimeInference & 4221 & 17015 & 58704.9500\\
RegionD & AITraining & 5726 & 407683 & 1387122.6833\\
RegionD & BatchInference & 3346 & 45291 & 155177.2000\\
RegionD & RealTimeInference & 488 & 1993 & 6558.5500\\
RegionE & AITraining & 4091 & 292000 & 978535.1167\\
RegionE & BatchInference & 2172 & 29548 & 100158.3333\\
RegionE & RealTimeInference & 649 & 2628 & 9054.7667\\
RegionF & AITraining & 4177 & 296143 & 1032935.0333\\
RegionF & BatchInference & 2033 & 27173 & 90263.0167\\
RegionF & RealTimeInference & 497 & 1996 & 6837.0667\\
\bottomrule
\end{tabular}
\end{table}

由表~\ref{tab:q1_group_statistics} 可见，各区域的任务组成存在明显差异，因而预测模型需要同时保留区域效应与任务类型效应。RegionD、RegionE 和 RegionF 的 AITraining 任务对应较大的 GPU 需求与 GPU-hour，而 RegionA、RegionB 和 RegionC 的 RealTimeInference 任务数相对较多。任务数、GPU 需求与 GPU-hour 所刻画的侧重点并不相同，因此调度容量不能仅依据任务数量判断，仍须结合持续时间和逐时重叠关系进行审计。

预测模型引入短期、日内和周内滞后项，并采用
\begin{equation}
z_t=\left[
\sin\frac{2\pi t}{24},
\cos\frac{2\pi t}{24},
\sin\frac{2\pi t}{168},
\cos\frac{2\pi t}{168}
\right]^{\mathsf T}
\end{equation}
表示周期特征。未截断预测量为
\begin{equation}
u_{rkt}=\beta_0+
\sum_{\ell\in\{1,2,3,24,48,168\}}\beta_\ell D'_{rk,t-\ell}
+\gamma^{\mathsf T}z_t+\alpha_r+\delta_k,
\end{equation}
最终预测结果取
\begin{equation}
\widehat D_{rkt}=\max\{0,u_{rkt}\}.
\end{equation}

% MMW-ID: CON-Q1-6
预测误差采用
\begin{align}
\mathrm{MAE}&=\frac{1}{n}\sum_j|y_j-\widehat y_j|,\\
\mathrm{RMSE}&=\sqrt{\frac{1}{n}\sum_j(y_j-\widehat y_j)^2}
\end{align}
评价，WAPE 仅在真实值绝对和为正时计算。最终测试的总体 MAE 为 $24.9733$ GPU，总体 RMSE 为 $35.6945$ GPU；测试区间形成 $432$ 条区域—任务类型逐小时预测记录，需求统计明细形成 $43200$ 行结果。

% MMW-LIMITATION: D2
方法契约原计划比较季节性基线并执行完整超参数网格，但实际运行采用固定的 $\epsilon=1.35$ 和 $\alpha=0.001$，未完成全部候选的验证选模。因此，现有误差仅反映已实现固定模型在保留测试区间上的表现。

\subsubsection{任务候选域与容量约束}

% MMW-ID: CON-Q1-1
任务 $i$ 的合法执行区域—开工时刻集合为
\begin{equation}
\Omega_i=\left\{(r,s):s\ge a_i,\ s+h_i\le F_i,\
\ell_{o_ir}\le\ell_i^{\max},\ \omega_{i,2406}(s)=0\right\}.
\end{equation}
若 $\Omega_i$ 为空，则任务及排除原因被记录，构造过程返回结构化状态。

% MMW-ID: CON-Q1-2
每个任务必须且只能选择一个候选，实时推理任务必须在到达时刻开工：
\begin{align}
\sum_{(r,s)\in\Omega_i}x_{irs}&=1,\\
s_i&=\sum_{(r,s)\in\Omega_i}sx_{irs},\qquad c_i=s_i+h_i,\\
s_i&=a_i,\qquad i\in\mathcal I_{\mathrm{RT}}.
\end{align}

% MMW-ID: CON-Q1-3
区域逐时 GPU、IT 功率和设施功率满足
\begin{align}
\sum_{i,s}g_i\chi_{it}(s)x_{irs}&\le G_r^{\max},\\
L_{rt}^{N}+\sum_{i,s}g_ip_{k_i}^{GPU}\chi_{it}(s)x_{irs}
&\le P_r^{IT,\max},\\
\pi_r\left(L_{rt}^{N}+\sum_{i,s}g_ip_{k_i}^{GPU}
\chi_{it}(s)x_{irs}\right)&\le P_r^{F,\max}.
\end{align}

% MMW-ID: CON-Q1-4
能源结算采用实际重叠时长：
\begin{align}
E_{rt}^{AI}&=\sum_{i,s}g_ip_{k_i}^{GPU}\omega_{it}(s)x_{irs},\\
L_{rt}^{AI,avg}&=\frac{E_{rt}^{AI}}{\Delta t},\\
L_{rt}^{F,avg}&=\pi_r(L_{rt}^{N}+L_{rt}^{AI,avg}).
\end{align}

% MMW-ID: CON-Q1-5
区域逐时峰值利用率和 GPU-hour 平均利用率分别为
\begin{align}
U_{rt}^{peak}
&=\frac{\sum_{i,s}g_i\chi_{it}(s)x_{irs}}{G_r^{\max}},\\
U_{rt}^{avg}
&=\frac{\sum_{i,s}g_i\omega_{it}(s)x_{irs}}
{G_r^{\max}\Delta t}.
\end{align}
两项指标仅在 $G_r^{\max}>0$ 时计算。最后 $24$ 小时的区域 GPU 利用情况见图~\ref{fig:gpu_utilization}。

\begin{figure}[htbp]
\centering
\includegraphics[width=0.88\textwidth]{fig_1_gpu_utilization.png}
\caption{最后 $24$ 小时各区域 GPU 利用率}
\label{fig:gpu_utilization}
\end{figure}

图~\ref{fig:gpu_utilization} 展示了统一时间轴上的区域资源占用变化，可用于识别任务集中到达形成的短时波动。峰值利用率反映任何正重叠任务造成的瞬时容量占用，而平均利用率反映按实际重叠时长折算的 GPU-hour，两者不能相互替代。最终容量合规性仍由逐任务区间重叠和瞬时 GPU、IT 与设施功率约束共同判定。

\subsubsection{实际调度算法与代表性结果}

任务按实时属性、有效截止、持续时间、GPU 需求和任务编号形成固定词典序，再在完整候选域中按等待时间、单向时延、加入任务后的峰值 GPU 利用率、区域及开工时刻依次尝试。每次放置后，GPU、IT 功率和设施功率约束均被重新审计。最后 $24$ 小时共有 $538$ 个实际到达任务获得执行区域及开工、完成时刻，其中 $68$ 个任务在第 $2400$ 至第 $2405$ 小时完成结清。表~\ref{tab:q1_schedule_examples} 给出了任务级代表性调度记录，而非完整甘特图。

\begin{table}[htbp]
\centering
\caption{最后 $24$ 小时任务级代表性调度记录}
\label{tab:q1_schedule_examples}
\begin{tabular}{rllllrrrrr}
\toprule
任务ID & 类型 & 来源 & 执行区 & 到达 & 开工 & 完成 & 等待 & 时延 & GPU\\
\midrule
3498 & RealTimeInference & RegionB & RegionB & 2376 & 2376 & 2377.1833 & 0 & 5.0 & 2.0\\
3809 & AITraining & RegionF & RegionF & 2376 & 2376 & 2377.6333 & 0 & 5.0 & 32.0\\
7403 & BatchInference & RegionD & RegionD & 2376 & 2376 & 2382.2333 & 0 & 5.0 & 10.0\\
7902 & BatchInference & RegionE & RegionE & 2376 & 2376 & 2376.4333 & 0 & 5.0 & 11.0\\
9267 & BatchInference & RegionE & RegionE & 2376 & 2376 & 2376.3667 & 0 & 5.0 & 7.0\\
9415 & AITraining & RegionD & RegionD & 2376 & 2376 & 2377.4667 & 0 & 5.0 & 119.0\\
11961 & AITraining & RegionE & RegionE & 2376 & 2376 & 2380.1000 & 0 & 5.0 & 111.0\\
14647 & BatchInference & RegionA & RegionA & 2376 & 2376 & 2377.2333 & 0 & 5.0 & 11.0\\
17820 & RealTimeInference & RegionC & RegionC & 2376 & 2376 & 2377.2500 & 0 & 5.0 & 1.0\\
18201 & BatchInference & RegionE & RegionE & 2376 & 2376 & 2376.6667 & 0 & 5.0 & 16.0\\
26377 & BatchInference & RegionB & RegionB & 2376 & 2376 & 2382.4333 & 0 & 5.0 & 17.0\\
29779 & BatchInference & RegionF & RegionF & 2376 & 2376 & 2381.6000 & 0 & 5.0 & 17.0\\
\bottomrule
\end{tabular}
\end{table}

表~\ref{tab:q1_schedule_examples} 中的记录均在第 $2376$ 小时到达并开工，但任务持续时间和 GPU 需求存在明显差异，说明小时级到达索引不能替代连续完成时刻。表中既包含实时推理，也包含批量推理和 AI 训练任务，可用于核对即时开工、连续占用和任务类型映射。该表仅用于展示任务级字段和时间口径，完整 $50000$ 行任务轨迹见 \texttt{q1\_q2\_q4\_task\_schedule.csv}，不得将该代表性样例解释为完整调度结果。

% MMW-ID: CON-Q1-7
理论模型规定候选耗尽时在局部冲突集中执行不超过 $N_{\mathrm{back}}=2000$ 个节点的有限回溯，并在搜索截止前仍未形成完整方案时返回 \texttt{feasibility\_unresolved}。实际实现未执行该回溯，只保留确定性贪心可行构造。

% MMW-LIMITATION: D1
实际任务调度未实现有限冲突回溯，其探索范围小于方法契约规定的范围。

% MMW-LIMITATION: D8
正式任务构造采用内部时间截止，并为情景复算、审计和文件输出预留运行时间；截止时未完成的部分方案不得作为正式结果。

% MMW-LIMITATION: D10
试跑阶段只读取真实输入并检查有限任务样本的合法时间窗，因此试跑通过不能替代正式全时域可行性审计。

\subsection{问题二：无储能碳感知任务调度}

% MMW-ID: OBJ-Q2
问题二从问题一的可行任务方案出发，在理论模型中以单任务迁移或开工时刻变更构造候选，并对成本、碳排放、有效网络时延、新能源未利用率和等待量进行稳健尺度标准化。实际程序未执行该局部改进，而是复算确定性任务方案。

% MMW-ID: CON-Q2-1
全部 $50000$ 个实际任务均执行唯一选址、实时即时开工、到达与截止、单向时延、瞬时 GPU、IT 功率、设施功率和真实重叠电量约束，预测需求不替代真实任务。

% MMW-ID: CON-Q2-2
问题二不配置储能，故
\begin{equation}
R_{rt}^{ch}=G_{rt}^{ch}=C_{rt}=D_{rt}=0,
\end{equation}
并满足
\begin{align}
R_{rt}^{use}+R_{rt}^{sell}+R_{rt}^{curt}&=R_{rt}^{av},\\
G_{rt}^{buy}+R_{rt}^{use}&=L_{rt}^{F,avg},\\
G_{rt}^{sell}&=R_{rt}^{sell}.
\end{align}

% MMW-ID: CON-Q2-3
购售电边界满足
\begin{align}
0\le G_{rt}^{buy}&\le I_r^{\max}y_{rt}^{grid},\\
0\le G_{rt}^{sell}&\le X_r^{grid}(1-y_{rt}^{grid}),
\qquad y_{rt}^{grid}\in\{0,1\}.
\end{align}
问题二不读取储能售电限制或 SOC 字段。

% MMW-ID: CON-Q2-4
新能源优先分流规则为
\begin{align}
R_{rt}^{use}&=\min\{R_{rt}^{av},L_{rt}^{F,avg}\},\\
G_{rt}^{buy}&=L_{rt}^{F,avg}-R_{rt}^{use}.
\end{align}
新能源富余时，外送量受 $X_r^{grid}$ 限制，其余部分被计为弃电。

% MMW-ID: CON-Q2-5
问题二运行成本为 $-419658147.7076$ 元，碳排放为 $14.3916$ tCO$_2$。全时域累计购电为 $40.4664$ MWh，累计售电为 $1329591.2956$ MWh。由于成本量纲为购电支出减售电收入，负值表示在给定售电价格及外送边界下累计售电收入超过购电支出，并不表示发电成本为负。

% MMW-ID: CON-Q2-6
有效网络时延为 $5.1594$ ms，新能源利用率为 $0.679078$。任务方案中有 $209$ 个任务的执行区域不同于来源区域，末端结清任务为 $68$ 个；GPU、IT 功率、设施功率、网络时延和完成时限审计均通过。

% MMW-LIMITATION: D3
实际程序未执行契约规定的单任务局部改进，$209$ 个迁移任务来自确定性贪心方案，而不是局部搜索收敛结果。

\subsection{问题三：固定负荷下的储能协同模型}

% MMW-ID: OBJ-Q3
问题三固定采用基准 AI IT 负荷和非 AI IT 负荷，并在理论模型中评价成本、碳排放、非负峰值净购电、绝对爬坡及有效新能源未利用率。无储能状态只作为比较对象，不作为含储能构造的启动判据。

% MMW-ID: CON-Q3-1
区域设施负荷及能源流满足
\begin{align}
L_{rt}^{F,avg}&=\pi_r(L_{rt}^{BAI}+L_{rt}^{N}),\\
R_{rt}^{use}+R_{rt}^{ch}+R_{rt}^{sell}+R_{rt}^{curt}&=R_{rt}^{av},\\
G_{rt}^{load}+R_{rt}^{use}+D_{rt}&=L_{rt}^{F,avg},\\
C_{rt}&=R_{rt}^{ch}+G_{rt}^{ch},\\
G_{rt}^{buy}&=G_{rt}^{load}+G_{rt}^{ch},\qquad
G_{rt}^{sell}=R_{rt}^{sell}.
\end{align}

% MMW-ID: CON-Q3-2
储能状态从第 $0$ 小时运行前的初始状态递推：
\begin{align}
S_{r,-1}&=S_r^0,\\
S_{rt}&=S_{r,t-1}+\eta_r^cC_{rt}\Delta t
-\frac{D_{rt}\Delta t}{\eta_r^d},\\
S_r^{min}&\le S_{rt}\le S_r^{max},\\
S_{r,2406}&\ge S_r^0.
\end{align}

% MMW-ID: CON-Q3-3
充放电功率及购售电满足
\begin{align}
0\le C_{rt}&\le C_r^{\max},&
0\le D_{rt}&\le D_r^{\max},\\
0\le G_{rt}^{buy}&\le I_r^{\max}y_{rt}^{grid},&
0\le G_{rt}^{sell}
&\le\min\{X_r^{grid},X_r^{storage}\}(1-y_{rt}^{grid}).
\end{align}
储能充放电及电网购售电均被设置为互斥行为。

% MMW-ID: CON-Q3-4
无储能状态按 $C_{rt}=D_{rt}=0$ 构造。若其违反进口上限，则只记录 \texttt{no\_storage\_baseline\_infeasible}，含储能方案仍从 $S_{r,-1}=S_r^0$ 独立前推。

% MMW-ID: CON-Q3-5
每次储能动作后均检查终端 SOC 的必要可达性：
\begin{equation}
S_{rt}+\sum_{u=t+1}^{2406}\eta_r^cC_r^{\max}\Delta t\ge S_r^0.
\end{equation}
完整轨迹随后接受购电上限、能源守恒、SOC、互斥关系和终端状态审计。

% MMW-ID: CON-Q3-6
净购电、峰值及绝对爬坡量定义为
\begin{align}
N_{rt}&=G_{rt}^{buy}-G_{rt}^{sell},\\
P_r^{peak}&=\max\{0,\max_tN_{rt}\},\\
V_r&=\sum_{t=1}^{2406}|N_{rt}-N_{r,t-1}|.
\end{align}

% MMW-ID: CON-Q3-7
新能源利用率仅在累计可用新能源为正时计算，否则返回结构化不可用状态并退出候选评分。

实际程序采用规则型 SOC 前向策略。无储能与含储能运行成本分别为 $-419820327.7393$ 元和 $-419516103.8102$ 元，成本增加 $304223.9291$ 元。初始 SOC 合计为 $1594.0000$ MWh，终端 SOC 合计为 $3540.0000$ MWh；累计充电量为 $2076.3155$ MWh，累计放电量为 $0$ MWh，且电网充电量为 $0$。无储能与含储能状态的累计购电量均为 $31.5979$ MWh，而累计售电量分别为 $1330267.7686$ MWh 和 $1328783.5422$ MWh。因此，成本增加来自原可售新能源被转化为储能库存后造成的售电收入减少；购电量保持不变，故两种状态的碳排放均为 $8.8506$ tCO$_2$。相应 SOC 轨迹见图~\ref{fig:storage_soc}。

\begin{figure}[htbp]
\centering
\includegraphics[width=0.88\textwidth]{fig_2_storage_soc.png}
\caption{各区域储能 SOC 轨迹}
\label{fig:storage_soc}
\end{figure}

图~\ref{fig:storage_soc} 展示了规则型策略下储能状态的逐时变化，轨迹应结合初始 SOC、上下限和终端约束判断。累计放电量为 $0$ MWh，说明该规则没有形成依靠放电削峰或电价套利的循环。六区域非负峰值净购电功率之和在储能前后均为 $31.5979$ MW，绝对爬坡量均为 $63.1957$ MW，因而该候选未产生削峰或平滑收益。

% MMW-LIMITATION: D4
问题三未执行契约规定的储能动作候选搜索，成本增加及碳排放、峰值和爬坡量不变只描述当前规则型候选。

\subsection{问题四：算—储—电联合调度与情景复算}

% MMW-ID: OBJ-Q4
联合模型评价运行成本、碳排放、有效网络时延、有效等待损失、新能源未利用率及六区域非负峰值净购电功率。各维度仅在有定义、数值有限且候选之间存在正尺度时进入代理评分。

% MMW-ID: CON-Q4-1
联合设施负荷由真实任务方案重新构造：
\begin{equation}
L_{rt}^{F,avg}=\pi_r\left[
L_{rt}^{N}+\frac{1}{\Delta t}
\sum_{i,s}g_ip_{k_i}^{GPU}\omega_{it}(s)x_{irs}
\right].
\end{equation}
任务侧执行问题一的全部硬约束，能源侧执行问题三的能源守恒、SOC、功率和互斥约束。

% MMW-ID: CON-Q4-2
联合初始方案从 $S_{r,-1}=S_r^0$ 独立前推至第 $2406$ 小时，不因问题三无储能比较状态的可行性而停止；若声明预算内未形成完整候选，则返回 \texttt{feasibility\_unresolved}。

% MMW-ID: CON-Q4-3
实时即时开工率、SLA 满足率、按期完成率和等待质量均先执行分母检查。实际程序采用
\begin{equation}
Q_{\mathrm{service}}=\frac{1}{1+\overline W}
\end{equation}
作为平均等待小时的有限代理，结果为 $1.0$。

% MMW-LIMITATION: D7
服务质量未采用按最大允许等待时间归一化的完整指标，故该结果只能解释为平均等待代理。

% MMW-ID: CON-Q4-4
碳约束、峰谷价差、售电价格和新能源水平按方法契约中的离散集合构造，新能源波动情景采用固定种子 $2026$ 并执行非负截断。

% MMW-ID: CON-Q4-5
理论模型规定联合基准与全部情景共享候选评价上限，各情景从同一联合基线独立重构；实际运行共检查 $15$ 个规则型候选。

% MMW-ID: CON-Q4-6
所有候选在启动前和可中断循环内部检查单调时钟，达到搜索截止后丢弃尚未完成的候选，剩余时间只用于文件输出、完整审计和状态整理。

问题四的全部情景结果见表~\ref{tab:q4_scenarios}。

\begin{table}[htbp]
\centering
\caption{问题四的十五个情景比较结果}
\label{tab:q4_scenarios}
\begin{tabular}{llrrrrrrl}
\toprule
情景 & 水平 & 成本/元 & 碳排放 & 时延 & 服务质量 & 新能源利用率 & 峰值/MW & 状态\\
\midrule
基准 & 1.00 & -419353923.7785 & 14.3916 & 5.1594 & 1.0 & 0.679129 & 40.4664 & 完成\\
碳约束 & 0.00 & -419353923.7785 & 14.3916 & 5.1594 & 1.0 & 0.679129 & 40.4664 & 完成\\
碳约束 & 0.10 & -419353923.7785 & 14.3916 & 5.1594 & 1.0 & 0.679129 & 40.4664 & 启发式内未达到\\
碳约束 & 0.20 & -419353923.7785 & 14.3916 & 5.1594 & 1.0 & 0.679129 & 40.4664 & 启发式内未达到\\
碳约束 & 0.30 & -419353923.7785 & 14.3916 & 5.1594 & 1.0 & 0.679129 & 40.4664 & 启发式内未达到\\
峰谷价差 & 0.75 & -419352397.8360 & 14.3916 & 5.1594 & 1.0 & 0.679129 & 40.4664 & 完成\\
峰谷价差 & 1.00 & -419353923.7785 & 14.3916 & 5.1594 & 1.0 & 0.679129 & 40.4664 & 完成\\
峰谷价差 & 1.25 & -419355449.7210 & 14.3916 & 5.1594 & 1.0 & 0.679129 & 40.4664 & 完成\\
售电机制 & 0.00 & 12327.7690 & 14.3916 & 5.1594 & 1.0 & 0.679129 & 40.4664 & 完成\\
售电机制 & 0.50 & -209670798.0048 & 14.3916 & 5.1594 & 1.0 & 0.679129 & 40.4664 & 完成\\
售电机制 & 1.00 & -419353923.7785 & 14.3916 & 5.1594 & 1.0 & 0.679129 & 40.4664 & 完成\\
新能源水平 & 0.80 & -302537839.3757 & 30759.1226 & 5.1594 & 1.0 & 0.805166 & 597.8863 & 完成\\
新能源水平 & 1.00 & -419353923.7785 & 14.3916 & 5.1594 & 1.0 & 0.679129 & 40.4664 & 完成\\
新能源水平 & 1.20 & -457183971.0773 & 0.0000 & 5.1594 & 1.0 & 0.577105 & 0.0000 & 完成\\
新能源波动 & 2026 & -402529632.4199 & 7923.8416 & 5.1594 & 1.0 & 0.675122 & 935.7721 & 完成\\
\bottomrule
\end{tabular}
\end{table}

表~\ref{tab:q4_scenarios} 表明，所有情景均保持 $209$ 个迁移任务，网络时延和服务质量代理也保持不变，这与情景复算固定联合基准任务方案的实现方式一致。售电价格系数为 $0$ 时运行成本为正，而系数为 $0.5$ 和 $1$ 时运行成本为负，说明当前成本指标对售电收入核算较为敏感。新能源水平系数由 $0.8$ 变为 $1.2$ 时，碳排放分别表现为 $30759.1226$ tCO$_2$ 和 $0$ tCO$_2$，峰值净购电分别为 $597.8863$ MW 和 $0$ MW；新能源波动情景则得到 $7923.8416$ tCO$_2$ 的碳排放和 $935.7721$ MW 的峰值净购电。碳削减水平为 $0.1$、$0.2$ 和 $0.3$ 时，目标均在声明启发式内未达到，该状态不构成数学不可行性的证明。

% MMW-LIMITATION: D5
问题四未执行契约规定的共享交替邻域搜索，每个情景只重建并审计一个规则型候选。

% MMW-LIMITATION: D6
碳约束情景未执行额外的碳定向搜索，未达标状态只能表述为在声明启发式内未达到。

% MMW-LIMITATION: D9
只有新能源波动情景采用固定种子 $2026$ 生成随机扰动，其余情景均为确定性复算。

联合基准运行成本为 $-419353923.7785$ 元，碳排放为 $14.3916$ tCO$_2$，网络时延为 $5.1594$ ms，新能源利用率为 $0.679129$，六区域非负峰值净购电功率之和为 $40.4664$ MW。统一约束审计通过，结果被标记为 \texttt{full\_horizon\_feasible=true}。同时，结果被标记为 \texttt{non\_Pareto=true}，且不存在全局或局部最优性证书，因此只能被视为已声明规则及已检查情景中的可行候选。

\subsection{统一约束审计}

任务合法性按
\begin{align}
a_i&\le s_i,\qquad c_i\le F_i,\\
\ell_{o_ir_i}&\le\ell_i^{\max},\\
i\in\mathcal I_{\mathrm{RT}}&\Rightarrow s_i=a_i,\\
\omega_{i,2406}(s_i)&=0
\end{align}
逐项检查。问题二的能源残差以及问题三、四的能源与 SOC 残差分别定义为
\begin{align}
e_{rt}^{q2}
&=G_{rt}^{buy}+R_{rt}^{av}
-L_{rt}^{F,avg}-G_{rt}^{sell}-R_{rt}^{curt},\\
e_{rt}^{energy}
&=G_{rt}^{buy}+R_{rt}^{av}+D_{rt}
-L_{rt}^{F,avg}-C_{rt}-G_{rt}^{sell}-R_{rt}^{curt},\\
e_{rt}^{SOC}
&=S_{rt}-S_{r,t-1}-\eta_r^cC_{rt}\Delta t
+\frac{D_{rt}\Delta t}{\eta_r^d}.
\end{align}
只有任务唯一执行、实时开工、截止、时延、GPU、IT、设施功率、能源守恒、SOC、终端 SOC、购售电互斥、充放电互斥、外送上限及第 $2406$ 小时禁占用全部通过时，才输出全时域可行状态。

% MMW-LIMITATION: L1
实际方法未提供全局或局部最优性证书，所有结论均限于已构造并通过审计的启发式候选。

% MMW-LIMITATION: L2
任务候选和储能动作的完整局部邻域未被穷举，未检查邻域中可能存在更优可行方案。

% MMW-LIMITATION: L3
碳目标未达到只表示在声明启发式内未达到，不构成数学不可行性的证明。

% MMW-LIMITATION: L4
问题四各情景保持联合基准的任务迁移方案不变，无法反映任务迁移方向对情景变化的全部响应。

% MMW-LIMITATION: L5
储能策略未计入循环退化成本，工程经济性评价仍不完整。

% MMW-LIMITATION: L6
复杂实例可能在内部任务构造截止前返回 \texttt{feasibility\_unresolved}，该状态表示运行预算内未完成构造。

% MMW-LIMITATION: L7
问题二只使用区域级最大购电与最大外送功率，不读取储能售电限制或 SOC 字段。

% MMW-LIMITATION: L8
试跑文件仅证明真实输入基础门禁和有限样本时间窗检查通过，不能替代正式全时域审计。

\section{模型的检验与灵敏度分析}

模型检验由完整约束审计和参数灵敏度实验共同完成。前者检查任务唯一执行、实时任务即时开工、到达与截止、单向时延、GPU 容量、IT 功率、设施功率、第 $2406$ 小时禁占用、能源守恒、SOC 递推、终端 SOC、充放电互斥和购售电边界。实际运行的约束满足状态为真，且全时域可行标记为真，说明输出轨迹在所采用数值容差和附件边界下保持内部一致。该检验属于代理验证和物理一致性检验，不使用不可观测的最优标签。

灵敏度实验以问题二系统运行成本 $-419658147.7075976$ 元为基准，分别改变购电价格水平和新能源出力水平。实验直接复算扰动后的目标值及相对变化，不从曲线外推未测试的参数区间。

\subsection{购电价格灵敏度}

购电价格系数变化下的运行成本见表~\ref{tab:price_sensitivity}，对应曲线见图~\ref{fig:sensitivity_price}。

\begin{table}[htbp]
\centering
\caption{购电价格变化对系统运行成本的影响}
\label{tab:price_sensitivity}
\begin{tabular}{rrrr}
\toprule
参数变化 & 参数系数 & 运行成本（元） & 相对变化（\%）\\
\midrule
$-20\%$ & $0.8$ & $-419660613.261396$ & $-0.0005875148169636924$\\
$-10\%$ & $0.9$ & $-419659380.48449683$ & $-0.00029375740848894776$\\
$10\%$ & $1.1$ & $-419656914.9306984$ & $0.00029375740848894776$\\
$20\%$ & $1.2$ & $-419655682.15379924$ & $0.0005875148169636924$\\
\bottomrule
\end{tabular}
\end{table}

表~\ref{tab:price_sensitivity} 显示，购电价格下降 $20\%$ 时，运行成本相对基准变化 $-0.0005875148169636924\%$；价格上升 $20\%$ 时，相对变化为 $0.0005875148169636924\%$。当扰动缩小至 $10\%$ 时，目标变化的绝对值相应为 $0.00029375740848894776\%$。因此，在本次购售电结构和固定调度方案下，目标函数对整体购电价格水平表现出较强稳健性。

\begin{figure}[htbp]
\centering
\includegraphics[width=0.82\textwidth]{sensitivity_electricity_price.png}
\caption{系统运行成本对购电价格水平的灵敏度}
\label{fig:sensitivity_price}
\end{figure}

图~\ref{fig:sensitivity_price} 呈现了价格系数变化与运行成本之间近似对称的响应方向，价格提高使目标值上升，价格降低使目标值下降。曲线在已测试区间内的纵向变化相对基准成本较小，与表~\ref{tab:price_sensitivity} 中较低的相对变化幅度一致。该稳健性只适用于整体价格同比例缩放，不能直接推广至峰谷价差、售电价格或任务调度随价格联动变化的情形。

\subsection{新能源水平灵敏度}

新能源出力水平变化下的运行成本见表~\ref{tab:renewable_sensitivity}，对应曲线见图~\ref{fig:sensitivity_renewable}。

\begin{table}[htbp]
\centering
\caption{新能源水平变化对系统运行成本的影响}
\label{tab:renewable_sensitivity}
\begin{tabular}{rrrr}
\toprule
参数变化 & 参数系数 & 运行成本（元） & 相对变化（\%）\\
\midrule
$-20\%$ & $0.8$ & $-302819505.6070327$ & $27.84138536063257$\\
$-10\%$ & $0.9$ & $-378911868.9949424$ & $9.709397740812053$\\
$10\%$ & $1.1$ & $-445329416.4918755$ & $-6.1171858391189184$\\
$20\%$ & $1.2$ & $-457431461.6153355$ & $-9.000972366216745$\\
\bottomrule
\end{tabular}
\end{table}

新能源水平下降 $20\%$ 时，运行成本相对基准变化 $27.84138536063257\%$；下降 $10\%$ 时，相对变化仍达到 $9.709397740812053\%$。新能源水平上升 $10\%$ 和 $20\%$ 时，运行成本相对变化分别为 $-6.1171858391189184\%$ 和 $-9.000972366216745\%$。由此可得，模型目标对新能源水平较为敏感，且新能源下降造成的成本不利变化在已测试区间内更为显著。

\begin{figure}[htbp]
\centering
\includegraphics[width=0.82\textwidth]{sensitivity_renewable_level.png}
\caption{系统运行成本对新能源出力水平的灵敏度}
\label{fig:sensitivity_renewable}
\end{figure}

图~\ref{fig:sensitivity_renewable} 表明，随着新能源水平提高，运行成本持续下降，说明新能源直接消纳与外送收益对目标具有重要影响。曲线响应并不呈现与价格实验相同的小幅对称变化，反映了新能源不足时购电缺口扩大与新能源富余时外送受限之间的非线性差异。因而，在模型推广和运行决策中，新能源预测误差应被视为比整体购电价格同比例误差更重要的不确定性来源。

综合两组实验可知，系统运行成本对购电价格整体缩放表现稳健，而对新能源水平变化较为敏感。该结论建立在任务方案和能源分流规则保持一致的条件下；由于实际实现没有针对每个扰动重新执行完整任务邻域与储能动作搜索，灵敏度结果衡量的是固定启发式结构下的参数响应，而不是重新优化后的价值函数灵敏度。

\section{模型的评价与推广}

\subsection{模型的优点}

本文模型实现了任务瞬时资源占用与小时能源结算的分离。任务只要与某小时存在正重叠，其 GPU 和瞬时 IT 功率便被计入容量审计，而实际重叠时长被用于计算 GPU-hour 和电量。该处理能够同时适应分钟级任务和小时级电力数据，避免整小时计费造成的能源高估，也避免小时平均负荷掩盖瞬时容量冲突。

算力、设施与能源边界被置于统一可审计框架中。AI IT 负荷由任务类型功率映射和真实调度轨迹重建，非 AI IT 负荷保持不可迁移，PUE 作用于总 IT 负荷；新能源则被拆分为直接消纳、储能充电、外送和弃电。由此，成本、碳排放和新能源利用率均可依据任务与能源轨迹重新计算，而不依赖附件中的基准结果列。

模型还明确处理了时域末端和储能初始状态。第 $2400$ 至第 $2405$ 小时只承接此前到达的弹性任务，第 $2406$ 小时只用于终端结算；储能状态从第 $0$ 小时运行前的绝对初始 SOC 开始递推，并满足终端 SOC 约束。无储能比较状态与含储能构造保持逻辑独立，避免前者的可行性状态被错误传播至后者。

指标分母门禁和完整残差审计进一步提高了结果的可解释性。网络时延、新能源利用率和服务质量仅在相应分母有效时计算，不可用维度不会通过任意小常数强行进入评分。任务唯一执行、容量、功率、时延、截止时间、能源守恒、SOC 和互斥关系均在完整时域内接受复核。

跨情景结果揭示了不同外部因素的影响方向。售电价格系数为 $0$、$0.5$ 和 $1$ 时，运行成本分别为 $12327.7690$ 元、$-209670798.0048$ 元和 $-419353923.7785$ 元，而碳排放、网络时延、服务质量代理、新能源利用率和峰值净购电保持不变，表明该组情景的成本变化主要由售电收入核算驱动。峰谷价差系数为 $0.75$、$1$ 和 $1.25$ 时，成本分别为 $-419352397.8360$ 元、$-419353923.7785$ 元和 $-419355449.7210$ 元，其余报告指标保持不变。

新能源情景呈现出更显著的能源侧变化。新能源水平系数为 $0.8$ 时，成本、碳排放和峰值净购电分别为 $-302537839.3757$ 元、$30759.1226$ tCO$_2$ 和 $597.8863$ MW；系数为 $1.2$ 时，相应结果为 $-457183971.0773$ 元、$0$ tCO$_2$ 和 $0$ MW。固定种子 $2026$ 的新能源波动情景得到 $-402529632.4199$ 元的成本、$7923.8416$ tCO$_2$ 的碳排放和 $935.7721$ MW 的峰值净购电，说明新能源供给水平及波动会同时影响购售电收益、碳排放和电网峰值。

\subsection{模型的不足}

模型的主要不足在于数学模型与实际实现之间仍存在差距。需求预测采用固定 Huber 超参数，没有完成全部季节性候选及超参数组合的验证选模；任务调度没有执行契约规定的有限冲突回溯，问题二没有执行单任务局部改进，问题三没有执行储能动作候选搜索，问题四也没有执行基准与情景共享的大规模交替邻域搜索。因此，所得结果只能说明已构造候选通过全时域硬约束审计，无法排除未检查邻域中存在成本、碳排放或峰值表现更优的可行方案。

储能策略的工程刻画和动作机制均较为简化。规则策略仅利用新能源富余充电，累计充电量为 $2076.3155$ MWh，累计放电量为 $0$ MWh，电网充电量亦为 $0$；初始 SOC 合计为 $1594.0000$ MWh，终端 SOC 合计为 $3540.0000$ MWh。由于没有放电，储能不能减少购电量、峰值净购电或绝对爬坡量，也没有形成电价套利循环。

问题三的成本增加可由能源流分解得到直接解释。无储能与含储能状态的累计购电量均为 $31.5979$ MWh，而累计售电量分别为 $1330267.7686$ MWh 和 $1328783.5422$ MWh；相应运行成本分别为 $-419820327.7393$ 元和 $-419516103.8102$ 元。当前规则把部分原可外送新能源转化为终端库存，减少了售电收入，因而使成本增加 $304223.9291$ 元；购电量未改变，故碳排放保持为 $8.8506$ tCO$_2$。该结果只能说明现行只充不放规则未实现经济性与削峰改善，不能用于否定经过优化的储能调度。

成本指标采用“购电支出减售电收入”的货币量纲。问题二累计购电为 $40.4664$ MWh，累计售电为 $1329591.2956$ MWh，运行成本为 $-419658147.7076$ 元。负值表示在附件给定售电价格和区域最大外送边界下，累计售电收入超过购电支出，并不表示发电成本、储能成本或系统物理能耗为负。由于当前模型未纳入新能源建设、运维及储能退化等成本，该指标更准确地反映既定核算边界内的净购售电现金流。

通信模型只采用确定的单向网络时延，没有考虑带宽拥塞、迁移数据量、传输能耗和通信费用。服务质量采用 $1/(1+\text{平均等待小时})$ 的有限代理，而非按最大允许等待时间归一化的完整指标。对于网络资源紧张或任务迁移代价显著的实际系统，这些省略可能改变任务迁移方向和多目标权衡关系。

情景分析仍受固定任务方案限制。十五个情景均保持 $209$ 个迁移任务，网络时延为 $5.1594$ ms，服务质量代理为 $1.0$，说明情景变化主要通过能源流和储能轨迹复算体现，尚未反映任务迁移对价格、碳强度和新能源变化的响应。碳削减水平为 $0.1$、$0.2$ 和 $0.3$ 时，目标均在声明启发式内未达到；由于没有执行碳定向搜索，这一状态不构成数学不可行性证明。新能源波动分析也只采用固定种子 $2026$ 的单个情景，后续仍需在保留可复现种子的条件下开展重复仿真。

所有实际结果均被标记为 \texttt{non\_Pareto=true}，且没有全局或局部最优性证书。现有方案仅是确定性或固定种子启发式生成的 scenario-feasible 候选，不能被解释为 Pareto 前沿、局部最优解或连续与离散决策空间中的最优方案。

\subsection{模型的推广}

在算力调度方面，候选域方法可推广至多数据中心、边缘节点与云中心协同的任务分配问题。通过加入节点可用性、数据驻留、可信执行环境或任务亲和性约束，可处理具有隐私和合规要求的跨区域计算。对于允许抢占或分布式训练的任务，可将单一开工变量扩展为逐时执行份额，并增加任务连续性、检查点和同步通信约束。

在能源调度方面，区域独立平衡可扩展为含输电线路的网络模型。通过引入节点相角、线路容量或直流潮流约束，可研究任务迁移与电力跨区传输之间的替代关系。储能模型还可加入循环退化、温度、自放电、寿命成本和备用容量，使短期购售电收益、削峰需求与长期设备损耗得到统一权衡。

在不确定性处理方面，新能源出力、任务到达和执行时长可被构造为情景树或分布鲁棒集合。实际运行时，可在每个小时更新任务队列、负荷预测和新能源预测，并保留已开工任务的连续占用状态，从而形成可在线实施的模型预测控制框架\cite{rawlings2017model}。对于新能源波动，应采用多个预先固定的随机种子重复复算，并分别报告能源、碳排放与峰值指标的稳定性。

在求解方法方面，可先以本文确定性贪心策略生成可行初始方案，再采用混合整数规划、大邻域搜索或分解协调方法改进任务与储能决策。若计算预算有限，可按时延、截止裕度和最低资源下界筛除任务候选；若需要多目标决策，则可采用 $\varepsilon$-约束法生成经过验证的非支配候选集\cite{miettinen1999nonlinear}。无论采用何种扩展，任务轨迹、能源轨迹与评价指标均应在完整时域上重新计算，并继续保留结构化不可用状态、终止状态和约束残差审计。

\appendix
\section{程序代码}
完整可运行程序见提交包中的 \texttt{code/solution.py}。

\bibliographystyle{plain}
\bibliography{references}

\end{document}
